\paragraph{Supplement to Finding A52: the differential of the actual point relation.}
Let $v:D_r\to D_s$ be a finite etale refinement of degree prime
to $p$ in the compatible pointed tower of A51. Let $P_r,Q_r$
be its compatible rational lifts of distinct points, and let
$m_r,m_s$ be the exact orders of the actual $p$-primary classes
$[Q_r-P_r]$ and $[Q_s-P_s]$. Suppose that $m_r>m_s$.
Choose the rational principal functions from A51, so that
\[
 \operatorname{div}(f_i)=m_i(Q_i-P_i),\qquad
 N_v(f_r)=c f_s^{m_r/m_s},\quad c\in k^\times.
\tag{PD.1}
\]
Then the actual point relation gives a nonzero regular differential
\[
 0\ne\omega_r=\frac{df_r}{f_r}
       \in H^0(D_r,\Omega^1_{D_r/k}),\qquad
                         \operatorname{Tr}_v(\omega_r)=0.
\tag{PD.2}
\]
This conclusion is about the given point divisor, not an abstract
torsion module. It detects a new trace-zero differential at a
strict order jump; it supplies no upper bound on $m_r$ or the
annihilator index in (IA.14).

\begin{proof}
The divisibility $m_s\mid m_r$ and the norm formula are (NC.4)--(NC.5),
using the divisor-of-norm identity
\cite[Lemma~42.18.1, Tag~\texttt{02RT}]{Stacks02RT}.
Since both orders are powers of $p$ and the inequality is strict,
$p\mid m_r$ and $p\mid m_r/m_s$.
At any geometric point write $f_r=t^a u$, with $t$ a uniformizer
and $u$ a unit. Every valuation $a$ is divisible by $p$:
it is $m_r$, $-m_r$, or zero. Thus
\[
                  df_r/f_r=a\,dt/t+du/u=du/u
\tag{PD.3}
\]
is regular. This proves regularity at every point.
The differential is nonzero. Over the perfect field, $df_r=0$
would make $f_r$ a $p$th power in the function field. Its root
would have divisor $(m_r/p)(Q_r-P_r)$, contradicting exact
order. The same argument proves separability of the function
$f_r:D_r\to\mathbf P^1$.

For a finite separable extension, differentiation of the product
of the conjugates proves the identity
\[
       \operatorname{Tr}_v(df_r/f_r)
                    =dN_v(f_r)/N_v(f_r).
\tag{PD.4}
\]
One may check it after passing to a Galois closure; both sides
descend. Here the trace on differentials is defined by the
function-field trace and the etale identification of relative
differentials, so it preserves regularity. Substitution of
(PD.1) into (PD.4) gives
$(m_r/m_s)\,df_s/f_s=0$ in characteristic $p$, proving (PD.2).

The direct ramification estimate for this particular function
also gives no upper bound. Its degree is $m_r$, and its zero
and pole fibres each consist of one point, of total ramification
index $m_r$. At each, the different exponent is at least $m_r$:
in a local parameter expression the derivative of the leading
term of order $m_r$ vanishes, while separability makes a later
derivative nonzero. Consequently Riemann--Hurwitz gives only
\[
 2g(D_r)-2=-2m_r+\deg\operatorname{Diff}(f_r),
              \qquad \deg\operatorname{Diff}(f_r)\ge2m_r
\tag{PD.5}
\]
from these two forced fibres
\cite[Lemma~53.12.4, Tag~\texttt{0C1F}]{StacksRH}.
Their lower contributions cancel the degree term; no upper
bound on $m_r$ follows from this estimate. The trace-zero
identity likewise retains the new differential rather than
forcing $f_r$ to be a $p$th power. These are limits of these
specific deductions, not a geometric realization of arbitrary
orders or a nonexistence theorem for a different point-based
argument.
\end{proof}
